\documentclass{article}
\usepackage{amsmath, amssymb, amsthm}
\usepackage{geometry}
\usepackage[CJKspace]{xeCJK}  % 한글 지원
\setmainfont{Times New Roman}
\setCJKmainfont{Malgun Gothic}  % 시스템에 따라 다른 한글 글꼴 사용 가능

\usepackage{tikz}
\usepackage{pdflscape}
\usetikzlibrary{trees}

\geometry{
 a4paper,
 left=10mm,
 right=10mm,
 top=15mm,
 bottom=15mm
 }
 
 \setlength{\parindent}{0pt}

\title{블랜더 쉐이딩 아이디어}
\author{사서림}
\date{2025.11.23.\text{\_\-\-\_}}

\begin{document}
\maketitle
\newpage

%\tableofcontents

\section*{Shading Idea 1: Discrete Tone Mapping with Matrix Distortion}

This model utilizes discrete thresholding for tone mapping and applies a color distortion matrix in specific intensity ranges to achieve a stylized look.

\[
\begin{aligned}
\mathbf{L}_c, \mathbf{L}_v &\in \mathbb{R}^3 \quad (\text{Light Direction, Light Color Vector}) \\
\mathbf{T}(u,v) &: [0,1]^2 \rightarrow \mathbb{R}^3 \quad (\text{Texture Color Map}) \\
\hat{n} &\in \mathbb{R}^3, \quad \|\hat{n}\|=1 \quad (\text{Surface Normal Vector}) \\
L &= \max(0, -\mathbf{L}_c \cdot \hat{n}) \quad (\text{Lambertian Intensity})
\end{aligned}
\]

The tone mapping function $h_f(L)$ is defined as:
\[
h_f(L) = 
\begin{cases} 
1 & L > x_1 \\
L & x_1 \ge L > x_2 \\
c_1 & x_2 \ge L > x_3 \\
0 & x_3 \ge L 
\end{cases}
\]

The color distortion vector function $\mathbf{v}(\mathbf{L}_v)$ is defined as:
\[
\mathbf{v}(\mathbf{L}_v) = 
\begin{cases} 
\mathbf{L}_v & L > x_1 \\
\mathbf{A} \mathbf{L}_v & x_1 \ge L > x_2 \quad (\text{Warm Distortion Matrix}) \\
\mathbf{L}_v & x_2 \ge L 
\end{cases}
\]

The final color is computed using the Hadamard product ($\odot$):
\[
\text{Color}_{\text{final}} = h_f(L) \cdot [ \mathbf{T}(u,v) \odot \mathbf{v}(\mathbf{L}_v) ]
\]

\noindent \textbf{Parameters:} 
$x_1, x_2, x_3 \in \mathbb{R}$ (Thresholds), $c_1 \in \mathbb{R}$ (Shadow Intensity), $\mathbf{A} \in \mathbb{R}^{3 \times 3}$ (Color Matrix).

\vspace{1cm}

\section*{Shading Idea 2.1: Continuous Gradient Matrix Mapping}

This model introduces continuous gradient matrices $\mathbf{A}(L)$ and $\mathbf{B}(L)$ at the boundaries of tone transitions to ensure smooth chromatic aberration or warm-tone shifts.

\[
h_f(L) = 
\begin{cases} 
1 & L > x_1 \\
c_1 & x_1 \ge L > x_2 \\
0 & x_2 \ge L 
\end{cases}
\]

\[
\mathbf{v}(\mathbf{L}_v) = 
\begin{cases} 
\mathbf{A}(L) \mathbf{L}_v & \delta_{1,1} \ge L > \delta_{1,2} \quad (\text{Warm/Dark Gradient Region}) \\
\mathbf{B}(L) \mathbf{L}_v & \delta_{2,1} \ge L > \delta_{2,2} \quad (\text{Warm/Bright Gradient Region}) \\
\mathbf{L}_v & \text{otherwise}
\end{cases}
\]

\noindent \textbf{Continuity Constraints:}
\[
\mathbf{A}(L), \mathbf{B}(L) \in C^k([0,1], \mathbb{R}^{3 \times 3}) \quad \text{where } k \ge 1
\]
\[
\mathbf{A}(\delta_{1,1}) = \mathbf{A}(\delta_{1,2}) = \mathbf{B}(\delta_{2,1}) = \mathbf{B}(\delta_{2,2}) = \mathbf{I}_{3 \times 3}
\]

\noindent \textbf{Parameter Valid Intervals:}
To ensure proper gradient transition regions between tones, the threshold parameters must satisfy the following inequality chain:
\[
1 \ge \delta_{1,1} > \delta_{1,2} > x_1 > \delta_{2,1} > \delta_{2,2} > x_2 > 0
\]
\begin{itemize}
    \item $[\delta_{1,2}, \delta_{1,1}]$: Gradient region before transitioning to the mid-tone ($x_1$).
    \item $[\delta_{2,2}, \delta_{2,1}]$: Gradient region before transitioning to the shadow tone ($x_2$).
\end{itemize}

\[
\text{Color}_{\text{final}} = h_f(L) \cdot [ \mathbf{T}(u,v) \odot \mathbf{v}(\mathbf{L}_v) ]
\]

\end{document}